% ============================================================
%  DIRECTED MOMENTUM FUNNELLING AND SELF-SELECTING COHERENCE
%  IN CONFINED COUPLED LATTICES
% ============================================================

\documentclass[twocolumn,10pt,a4paper]{article}

\usepackage[utf8]{inputenc}
\usepackage[T1]{fontenc}
\usepackage{lmodern}
\usepackage[a4paper,top=2cm,bottom=2.2cm,left=1.8cm,right=1.8cm,columnsep=0.6cm]{geometry}
\usepackage{amsmath,amssymb,amsthm,mathtools}
\usepackage{bm}
\usepackage{booktabs}
\usepackage{array}
\usepackage{enumitem}
\usepackage{graphicx}
\usepackage{hyperref}
\usepackage[round,sort&compress]{natbib}
\usepackage{microtype}
\usepackage{float}
\usepackage{xcolor}
\usepackage{fancyhdr}

\pagestyle{fancy}
\fancyhf{}
\fancyhead[LE,RO]{\thepage}
\fancyhead[RE]{\small Directed Momentum Funnelling in Confined Coupled Lattices}
\fancyhead[LO]{\small Sachikonye}
\renewcommand{\headrulewidth}{0.4pt}

\hypersetup{
  colorlinks=true,
  linkcolor=blue!50!black,
  citecolor=blue!50!black,
  urlcolor=blue!50!black,
  pdftitle={Directed Momentum Funnelling and Self-Selecting Coherence in Confined Coupled Lattices},
  pdfauthor={Kundai Farai Sachikonye}
}

\newtheoremstyle{thmstyle}{\topsep}{\topsep}{\itshape}{}{\bfseries}{.}{.5em}{}
\theoremstyle{thmstyle}
\newtheorem{theorem}{Theorem}[section]
\newtheorem{lemma}[theorem]{Lemma}
\newtheorem{proposition}[theorem]{Proposition}
\newtheorem{corollary}[theorem]{Corollary}
\theoremstyle{definition}
\newtheorem{definition}[theorem]{Definition}
\newtheorem{principle}[theorem]{Principle}
\theoremstyle{remark}
\newtheorem{remark}[theorem]{Remark}

\newcommand{\R}{\mathbb{R}}
\newcommand{\N}{\mathbb{N}}
\newcommand{\vsig}{v_{\mathrm{sig}}}
\newcommand{\vnode}{v_{\mathrm{el}}}
\newcommand{\xstar}{\mathbf{x}^{\ast}}

% ============================================================
\title{\textbf{Directed Momentum Funnelling and Self-Selecting\\
Coherence in Confined Coupled Lattices:\\
A Substrate-Independent Mechanism for\\
Single-Element State Delivery and Passive Recording}}

\author{
Kundai Farai Sachikonye\\
Independent Researcher\\
\texttt{kundai.sachikonye@bitspark.com}
}
\date{\today}
% ============================================================

\begin{document}
\maketitle

% ============================================================
\begin{abstract}
\noindent
We study a general class of physical systems --- \emph{confined coupled
lattices} --- consisting of identical elements held in a bounded region by
an external potential and interacting through a repulsive, inverse-square
coupling. We establish a sequence of structural properties of such systems
and show that, taken together, they constitute a mechanism for delivering a
single element to a high directed momentum while requiring no active
high-precision control. The mechanism rests on four results. \textbf{(I)
Funnelling disparity:} a perturbation propagates through the coupled lattice
at a signal velocity $\vsig=\sqrt{\kappa/(m\,d)}$ that diverges as the
inter-element spacing $d$ contracts and that exceeds the velocity of any
individual element by an arbitrary factor; consequently the momentum stored
across many elements can be delivered to one leading element over a time set
by $\vsig$, not by element transport. \textbf{(II) Compression cascade:} a
velocity-graded sequential release --- leading elements slowest --- forces
the lattice to self-compress, raising $\vsig$ monotonically along the chain
and funnelling momentum forward, in direct analogy with the tip of a tapered
elastic medium. \textbf{(III) Self-selecting coherence:} the coherent output
configuration is a fixed-point attractor with a basin of positive measure;
its existence is therefore its own confirmation, with no thresholding step
and a structurally vanishing false-positive rate. A stochastic source
operated at low per-trial success but high throughput yields output whose
purity is set by causal construction rather than by statistical subtraction.
\textbf{(IV) Passive recording:} a metastable variant of the same lattice
records the sequence of perturbing transits as persistent local
destabilizations, functioning as an integrating recorder that is written by
the transits themselves and read once. We further show that staged
confirmation chains produce single-outcome selectivity, that a threshold
conversion event is confirmed by detecting its reverse signature, and that
under a precise time reference the entire measurement reduces to timestamped
counting. The construction is substrate-independent; we close by noting the
regime in which it instantiates as a compact directed-delivery and recording
device.

\bigskip
\noindent\textbf{Keywords:} coupled lattices, signal velocity, momentum
funnelling, compression cascade, fixed-point selection, metastable
recording, confinement, self-organization, inverse-square coupling
\end{abstract}

\tableofcontents

% ============================================================
\section{Introduction}
\label{sec:intro}
% ============================================================

\subsection{The problem and an organizing principle}

The delivery of a single element of a many-body system to a high directed
momentum is conventionally approached by \emph{control}: one prepares a
coherent collective object and imposes precise external fields to drive it,
spending the engineering budget on suppressing every departure from the
intended trajectory. The cost of this approach scales with the precision
demanded and with the spatial extent over which the driving must be
maintained.

This paper develops an alternative in which the quantities that the
control-based approach suppresses are instead repurposed as the operating
mechanism. We organize the development around a single principle, which we
state once and apply repeatedly.

\begin{principle}[Inversion]
\label{prin:inversion}
In a confined coupled lattice, three phenomena ordinarily treated as
obstacles --- (i) the mutual repulsion that disperses a confined collection,
(ii) the instability of a driven collective state, and (iii) the high failure
rate of an uncontrolled release --- can each be configured as the mechanism
that performs, respectively, momentum delivery, recording, and selection.
\end{principle}

The repulsion that disperses a confined collection, when the collection is
released with an appropriate velocity grading, becomes the forward
momentum-funnelling interaction (Sec.~\ref{sec:cascade}). The instability of
a driven collective state, left uncorrected, becomes the write mechanism of a
passive recorder (Sec.~\ref{sec:recording}). The high failure rate of an
uncontrolled release, paired with a fixed-point selection criterion, becomes
a purity-by-construction filter (Secs.~\ref{sec:selection}--\ref{sec:chains}).
Each application is a consequence of structural properties of the lattice,
derived below.

\subsection{The system}

Throughout, the object of study is the following.

\begin{definition}[Confined coupled lattice]
\label{def:lattice}
A \emph{confined coupled lattice} is a system of $N$ identical elements of
mass $m$ and coordinates $\{\mathbf{x}_i\}_{i=1}^N$, evolving under the
Hamiltonian
\begin{equation}
H = \sum_{i=1}^N \frac{|\mathbf{p}_i|^2}{2m}
  + \sum_{i<j} \frac{\kappa}{|\mathbf{x}_i-\mathbf{x}_j|}
  + \sum_{i=1}^N V_{\mathrm{ext}}(\mathbf{x}_i),
\label{eq:hamiltonian}
\end{equation}
where $\kappa>0$ sets the strength of the repulsive inverse-square coupling
and $V_{\mathrm{ext}}$ is an external confining potential with bounded
sublevel sets.
\end{definition}

The pairwise potential $\kappa/r$ is the generic form of a like-signed
inverse-square coupling; systems realizing Eq.~\eqref{eq:hamiltonian} include
trapped single-species charged ensembles \citep{dubin1999,birkl1992},
charged colloidal suspensions, and the one-component plasma whose ground
state is the Wigner lattice \citep{wigner1934}. We make no further
commitment as to the physical identity of the elements; the results are
properties of Eq.~\eqref{eq:hamiltonian}.

\subsection{Overview of results}

Section~\ref{sec:equilibrium} establishes the equilibrium configuration: a
nested-shell lattice, with the linear chain as its one-dimensional limit.
Section~\ref{sec:release} formalizes stored collective modes and directed
release. Section~\ref{sec:signal} derives the signal velocity and the
funnelling disparity (Result I). Section~\ref{sec:cascade} derives the
compression cascade (Result II). Section~\ref{sec:transverse} establishes the
self-confinement of the co-moving lattice. Section~\ref{sec:selection}
develops self-selecting coherence (Result III), and Section~\ref{sec:chains}
extends it to staged confirmation chains. Section~\ref{sec:recording}
develops the metastable recorder (Result IV). Section~\ref{sec:threshold}
treats threshold conversion and reverse-signature confirmation.
Section~\ref{sec:counting} reduces measurement to timestamped counting.
Section~\ref{sec:discussion} discusses generality and the regime of
instantiation.

% ============================================================
\section{Equilibrium Configuration}
\label{sec:equilibrium}
% ============================================================

\subsection{Nested shells}

Minimizing Eq.~\eqref{eq:hamiltonian} under a cylindrically symmetric
confinement yields a configuration of concentric shells. This is the
generic ground-state morphology of confined inverse-square systems and is
observed directly in trapped charged ensembles, where laser cooling reveals
ordered concentric shells \citep{birkl1992,dubin1999}.

\begin{proposition}[Shell structure]
\label{prop:shells}
For $V_{\mathrm{ext}}$ harmonic and cylindrically symmetric, the minimizers
of Eq.~\eqref{eq:hamiltonian} organize into $K$ concentric shells, with the
intershell spacing approaching a constant set by the balance of confinement
and coupling, and the linear chain ($K=1$) recovered in the strong-confinement
limit transverse to the symmetry axis.
\end{proposition}

\noindent The proof is the standard variational analysis of one-component
confined plasmas \citep{dubin1999}; we use only the conclusion. The shell
index will serve below as a transverse coordinate, the concentric shells
providing mutual lateral support for one another
(Sec.~\ref{sec:transverse}).

\subsection{The chain limit and its instability}

The one-dimensional limit --- a single chain of elements along the symmetry
axis --- is the configuration in which the funnelling mechanism is cleanest,
because all coupling forces are longitudinal. It is, however, unstable to
transverse buckling: beyond a critical linear density the chain undergoes a
structural transition to a zig-zag and then to multiple shells
\citep{fishman2008}. We will see (Sec.~\ref{sec:transverse}) that the
nested-shell configuration of Proposition~\ref{prop:shells} is precisely the
stable endpoint of this transition, so that the multi-shell lattice realizes
many longitudinal chains held mutually rigid --- the chain limit made stable
by multiplicity.

% ============================================================
\section{Stored Collective Modes and Directed Release}
\label{sec:release}
% ============================================================

\subsection{Stored rotational energy}

Let the confined lattice be set into rigid collective rotation at angular
velocity $\omega$ about the symmetry axis. An element at radius $r$ then
carries tangential speed
\begin{equation}
v_{\mathrm{tan}}(r) = \omega r ,
\label{eq:tangential}
\end{equation}
and the stored kinetic energy of the rotating configuration is
$E_{\mathrm{rot}} = \tfrac12 I \omega^2$ with $I$ the moment of inertia about
the axis. The rotation is a collective mode: it stores energy in the
configuration without displacing any element from its shell
\citep{goldstein2002,arnold1989}.

\subsection{Directed release}

\begin{definition}[Directed release]
\label{def:release}
A \emph{directed release} is the removal of the confining constraint on a
single element at a chosen phase $\phi$ of the collective rotation. The
released element departs with the instantaneous tangential velocity
$v_{\mathrm{tan}}(r)$ of Eq.~\eqref{eq:tangential}, directed along the
tangent at phase $\phi$.
\end{definition}

Directed release converts stored collective-mode energy into the directed
translational energy of one element. Because only the released element leaves,
the collective mode is perturbed only at order $1/N$; the configuration acts
as a quasi-stationary reservoir from which single elements are drawn at the
rim speed. We emphasize that release at a definite phase $\phi$ requires only
that the constraint be opened, not that the element be steered: the direction
is supplied by the rotation, not by an external field.

% ============================================================
\section{Signal Velocity and the Funnelling Disparity}
\label{sec:signal}
% ============================================================

\subsection{Propagation of a longitudinal perturbation}

Consider a chain of elements at uniform spacing $d$ along the symmetry axis.
The longitudinal stiffness of a single bond is the curvature of the pair
potential at the equilibrium spacing,
\begin{equation}
k(d) = \left.\frac{d^2}{dr^2}\frac{\kappa}{r}\right|_{r=d}
     = \frac{2\kappa}{d^3}.
\label{eq:stiffness}
\end{equation}
A chain of masses $m$ coupled by springs of stiffness $k$ and spacing $d$
transmits longitudinal disturbances at the lattice sound speed
\citep{ashcroft1976,toda1967}
\begin{equation}
\vsig(d) = d\sqrt{\frac{k(d)}{m}}
         = \sqrt{\frac{2\kappa}{m\,d}}.
\label{eq:vsig}
\end{equation}

\begin{theorem}[Funnelling disparity]
\label{thm:disparity}
The signal velocity $\vsig(d)$ of Eq.~\eqref{eq:vsig} increases without bound
as $d\to 0$, scaling as $d^{-1/2}$, and is independent of the velocity of any
individual element. Hence a perturbation injected at one end of the chain
arrives at the far end after a time $L/\vsig$, where $L$ is the chain length,
which can be made arbitrarily small relative to the transit time $L/\vnode$
of any element moving at speed $\vnode$.
\end{theorem}

\begin{proof}
Equation~\eqref{eq:vsig} gives $\partial \vsig/\partial d<0$ and
$\vsig\to\infty$ as $d\to0$. The element velocity $\vnode$ enters neither
$k(d)$ nor $\vsig$; the disturbance is carried by the coupling, not by element
transport. The ratio $\vsig/\vnode$ is therefore unbounded as $d$ or $\vnode$
is varied independently.\qed
\end{proof}

\subsection{Consequence: momentum delivery without transport}

Theorem~\ref{thm:disparity} is the central enabling fact. A momentum stored
across many elements of the chain can be delivered to the leading element in
the time $L/\vsig$ taken by the perturbation to traverse the chain --- a time
governed by the coupling, while the elements themselves barely move. The
acceleration imparted to the leading element is then
\begin{equation}
a_{\mathrm{lead}} \sim \frac{\sum_i p_i}{m}\cdot\frac{\vsig}{L},
\label{eq:lead-accel}
\end{equation}
large precisely because $\vsig\gg\vnode$. The total deliverable momentum is
bounded above by the stored momentum $\sum_i p_i$ (no momentum is created);
the mechanism is one of \emph{funnelling}, concentrating distributed momentum
onto one element, not of generation.

\begin{remark}
The mechanism is the same as the transfer of a collision through a line of
near-touching equal masses, in which the disturbance traverses the line far
faster than any mass moves and a single mass departs the far end
\citep{hutzler2004}. Equation~\eqref{eq:vsig} is the continuum statement of
that transfer for an inverse-square-coupled chain.
\end{remark}

% ============================================================
\section{The Compression Cascade}
\label{sec:cascade}
% ============================================================

\subsection{Graded release}

Consider releasing a sequence of $n$ elements along the axis, the leading
element first, with speeds
\begin{equation}
v_1 < v_2 < \cdots < v_n ,
\label{eq:grading}
\end{equation}
so that each trailing element is faster than the one ahead of it. This is the
\emph{velocity-graded release}. Its defining consequence is geometric: the
trailing elements close on the leading ones, the chain self-compresses, and
the inter-element spacing $d$ decreases with time.

\begin{theorem}[Self-compression and forward funnelling]
\label{thm:cascade}
Under the grading of Eq.~\eqref{eq:grading}, the chain spacing $d(t)$ is
monotonically decreasing while the elements remain ordered. By
Eq.~\eqref{eq:vsig}, the signal velocity $\vsig(d(t))$ is then monotonically
increasing, and the momentum of the trailing elements is funnelled forward to
the leading element through the compressing chain.
\end{theorem}

\begin{proof}
Order the elements $1,\dots,n$ from front to back with positions
$x_1>x_2>\cdots>x_n$ and velocities satisfying Eq.~\eqref{eq:grading}. While
the ordering holds, $\dot x_i - \dot x_{i+1} = v_i - v_{i+1} <0$, so each gap
$x_i-x_{i+1}$ shrinks and the mean spacing $d(t)$ decreases. By
Eq.~\eqref{eq:vsig}, $\vsig\propto d^{-1/2}$ increases. The compression raises
the bond stiffness $k(d)\propto d^{-3}$ (Eq.~\eqref{eq:stiffness}), so each
successive bond transmits momentum forward more stiffly than the last; the
net longitudinal momentum flux is toward the leading element.\qed
\end{proof}

\subsection{The tapered-medium analogy}

Theorem~\ref{thm:cascade} reproduces, in a discrete setting, the
acceleration of the free end of a tapered elastic medium: a disturbance
launched into the thick end funnels its energy into the progressively lighter
tip, whose speed grows as the cross-section shrinks. Here the role of the
geometric taper is played by the velocity grading, which drives the
compression that raises $\vsig$ toward the front. The grading of
Eq.~\eqref{eq:grading} is the only ingredient required; release the same
elements in the opposite order and the chain disperses, $\vsig$ falls, and no
funnelling occurs. The sign of the grading is therefore the entire stability
criterion of the cascade.

\subsection{Recovery of the dense-lattice limit}

As compression proceeds, $d$ approaches the equilibrium spacing of a dense
lattice and $\vsig$ approaches the intrinsic sound speed of that lattice. The
graded-release chain thus interpolates continuously between a sparse,
slowly-propagating configuration and a dense, fast-propagating one; the
fully-compressed limit is an ordinary dense coupled lattice
\citep{ashcroft1976}. The cascade is the controlled traversal of this
interpolation.

% ============================================================
\section{Self-Confinement of the Co-Moving Lattice}
\label{sec:transverse}
% ============================================================

The chain limit is transversely unstable (Sec.~\ref{sec:equilibrium}); we now
show that the co-moving multi-shell lattice is not, and that its transverse
stability \emph{improves} with the very speed the cascade produces.

\subsection{Two opposing transverse forces}

A set of co-moving like-coupled chains constitutes parallel co-directed
currents of the coupling charge. Such currents exert two transverse forces on
one another: a static repulsion from the inverse-square coupling, and an
attraction from the motion-induced field of co-directed currents
\citep{bennett1934}. For elements co-moving at speed $v$, the attraction
grows with $v$, and the net transverse force per element scales as
\begin{equation}
F_\perp \;\propto\; \left(1-\beta^2\right) = \frac{1}{\gamma^2},
\qquad \beta=\frac{v}{c_{\ast}},\quad \gamma=(1-\beta^2)^{-1/2},
\label{eq:net-transverse}
\end{equation}
where $c_{\ast}$ is the propagation-speed ceiling of the medium (the limiting
value of $\vsig$).

\begin{corollary}[Speed-induced self-confinement]
\label{cor:selfconfine}
The net transverse defocusing force of Eq.~\eqref{eq:net-transverse} vanishes
as $\beta\to1$. The co-moving lattice therefore becomes transversely
force-free in the high-speed limit: the configuration that the cascade drives
the lattice toward is the configuration in which it is most stable.
\end{corollary}

This is the second instance of Principle~\ref{prin:inversion}: the dispersive
repulsion that destabilizes a static chain is, for a co-moving chain,
progressively cancelled by the current attraction, leaving the external
confinement responsible only for the low-speed head of the process. The
multi-shell morphology of Proposition~\ref{prop:shells} supplies the residual
lateral support: each shell is held by its neighbours, and the
linear-to-zig-zag-to-shell sequence \citep{fishman2008} terminates in exactly
the mutually rigid configuration required.

\begin{remark}[The hard regime is the start, not the run]
By Corollary~\ref{cor:selfconfine}, the transverse confinement burden is
maximal at $\beta\to0$ and falls monotonically thereafter. The design effort
concentrates entirely at initiation; the high-speed stage stabilizes itself.
\end{remark}

% ============================================================
\section{Self-Selecting Coherence}
\label{sec:selection}
% ============================================================

We now address operation without precise control. The claim is that a
suitably configured cascade produces a coherent output only from initial
conditions lying in a basin of positive measure, and that the existence of
the coherent output is therefore its own confirmation.

\subsection{The coherent output as a fixed point}

Let $\mathcal{S}$ denote the space of release configurations (phase, grading,
shell alignment) and let $T:\mathcal{S}\to\mathcal{S}$ be the map advancing a
release to the asymptotic state of the cascade.

\begin{definition}[Coherent attractor]
\label{def:attractor}
A \emph{coherent output} is a fixed point $\xstar=T(\xstar)$ of the cascade
map that is attracting on a basin $\mathcal{B}\subseteq\mathcal{S}$ of
positive measure: there is a neighbourhood norm under which $T$ is a
contraction on $\mathcal{B}$, so that every release in $\mathcal{B}$ converges
to $\xstar$ \citep{banach1922}.
\end{definition}

\begin{proposition}[Detection equals existence]
\label{prop:detection}
If the cascade map admits a coherent attractor with basin $\mathcal{B}$ of
positive measure, then a release either lies in $\mathcal{B}$ and converges to
the coherent output $\xstar$, or lies outside $\mathcal{B}$ and does not
produce $\xstar$. The presence of the coherent output is thus a binary,
all-or-nothing event: its existence certifies that the release lay in
$\mathcal{B}$, with no intermediate or partial output to be thresholded.
\end{proposition}

\begin{proof}
By Definition~\ref{def:attractor}, convergence to $\xstar$ occurs for and only
for releases in $\mathcal{B}$. The complement produces trajectories that do
not approach $\xstar$. Since the output is the fixed point or its absence,
there is no continuum of partial outputs; the readout is the indicator
function of $\mathcal{B}$.\qed
\end{proof}

Proposition~\ref{prop:detection} requires that the cascade be a contraction
toward a fixed point rather than a continuous amplifier; this is the precise
condition under which ``did the coherent output appear'' is a well-posed
binary question. The basin measure $|\mathcal{B}|$ is the single quantity on
which both yield and resolution depend, and its positivity is the central
physical requirement of the entire construction.

\subsection{Throughput-dominated yield}

Operate the source by releasing at a fixed rate, uncorrelated with the
collective phase, so that the release phase is uniform on $[0,2\pi)$. Then
the per-release success probability is
\begin{equation}
P = \frac{|\mathcal{B}|}{2\pi},
\qquad
\text{Yield} = N_{\mathrm{rel}}\,P\,\Gamma,
\label{eq:yield}
\end{equation}
where $N_{\mathrm{rel}}$ is the number of releases and $\Gamma$ the cascade
gain. Because each release involves a single element, a failed release
perturbs nothing and carries no penalty; there is no collective object to
spoil. The yield is therefore maximized by increasing $N_{\mathrm{rel}}$
(throughput), not by increasing $P$ (precision).

\begin{remark}
This is the third instance of Principle~\ref{prin:inversion}. A high failure
rate is not corrected; it is rendered costless by single-element operation and
compensated by throughput. The output is selected by the basin
$\mathcal{B}$, not by suppression of failures.
\end{remark}

% ============================================================
\section{Staged Confirmation Chains}
\label{sec:chains}
% ============================================================

The selection of Sec.~\ref{sec:selection} composes. Let stages
$k=1,\dots,K$ be cascades arranged so that the coherent output of stage $k$ is
the input release of stage $k+1$.

\begin{theorem}[Single-outcome selectivity]
\label{thm:chain}
In a staged confirmation chain, a completed run --- one producing coherent
output at the final stage --- occurs only when the coherent attractor is
reached at every stage. The output of the chain therefore contains only the
conjunction of all stage confirmations; any stage that fails to reach its
attractor terminates the run and produces no output.
\end{theorem}

\begin{proof}
By Proposition~\ref{prop:detection}, stage $k$ produces its coherent output
$\xstar_k$ iff its input lies in $\mathcal{B}_k$. Stage $k+1$ receives
$\xstar_k$ as input; absent $\xstar_k$ there is no input and stage $k+1$
cannot fire. By induction, the final output exists iff every stage reached its
attractor, i.e.\ iff the input to stage $1$ lay in the preimage
$\bigcap_k T_{1}^{-1}\cdots T_{k-1}^{-1}(\mathcal{B}_k)$. The output is thus
the conjunction of all stage confirmations.\qed
\end{proof}

\subsection{Purity by construction}

Theorem~\ref{thm:chain} has a consequence that distinguishes the chain from a
statistical filter. A configuration that fails any stage produces a
\emph{terminated} run --- no output. It is therefore not a contaminant in the
output to be subtracted; it is absent from the output entirely. The output
contains only completed chains, and a completed chain is, by construction, the
unique configuration consistent with every stage confirmation. Selectivity is
achieved causally, not by background subtraction.

\subsection{The completion-rate trade}

The completion probability of a $K$-stage chain is $P_{\mathrm{cplt}}
=\prod_k P_k$, which may be small. This is the deliberate trade: the chain
exchanges rate for purity. Because each stage after the first receives a
pre-conditioned input --- the coherent output of the previous stage, already
within tolerance --- the conditional pass probabilities $P_{k\geq2}$ approach
unity, and the selection cost is paid essentially once, at stage~$1$. The
chain delivers a low rate of unambiguous outputs in place of a high rate of
ambiguous ones, and the low rate is affordable for the same reason failures
are costless: single-element operation.

% ============================================================
\section{Passive Recording in a Metastable Lattice}
\label{sec:recording}
% ============================================================

We now use the lattice not as a source but as a recorder, by configuring it in
a metastable rather than a stable state and declining to correct the
perturbations imposed on it.

\subsection{The metastable configuration}

\begin{definition}[Metastable phase-locked lattice]
\label{def:metastable}
A \emph{metastable phase-locked lattice} is a confined coupled lattice held in
a locally stable but globally non-minimal configuration, such that a localized
perturbation nucleates a persistent local destabilization (a \emph{mark}) of
finite lifetime $\tau_{\mathrm{mark}}$, without triggering relaxation of the
whole lattice.
\end{definition}

This is the configuration of a system poised at a metastable point, where a
local disturbance nucleates a persistent feature rather than decaying or
propagating globally \citep{langer1969}. The requirement is a window: the
lattice must be stable enough to hold a mark (the mark does not self-erase)
yet metastable enough that the mark persists without correction (the lattice
does not relax it away).

\subsection{Transits write; the lattice integrates}

Let a sequence of external transits pass through the lattice, each
destabilizing the local region it crosses. In the metastable configuration the
destabilization persists; successive transits write marks into distinct
regions. The accumulated pattern of marks is the integrated record of the
transit sequence.

\begin{proposition}[Recording by accumulation]
\label{prop:recording}
In a metastable phase-locked lattice with mark lifetime $\tau_{\mathrm{mark}}$
and characteristic local relaxation (lag) time $\tau$, the lattice records a
transit sequence as a spatial pattern of persistent marks, with per-region
rate capacity $\tau^{-1}$ and total capacity set by the number of
distinguishable lattice regions. The record is read once, after accumulation.
\end{proposition}

\begin{proof}
By Definition~\ref{def:metastable}, each transit nucleates a mark that
persists for $\tau_{\mathrm{mark}}$. Two transits through the same region
within the local lag time $\tau$ are not resolved; hence the per-region rate
capacity is $\tau^{-1}$. Distinct regions record independently, giving total
capacity equal to the count of resolvable regions. Persistence over
$\tau_{\mathrm{mark}}$ permits deferred, single readout.\qed
\end{proof}

\begin{remark}[The recorder is written by what it measures]
The transits are simultaneously the objects measured and the styli that write
the record. No active element is energized during accumulation; the recorder
is a passive integrating medium, in the manner of historical track-recording
media in which a traversing object nucleates a persistent local feature in a
metastable bulk. This is the second application of
Principle~\ref{prin:inversion}: the instability that would have to be actively
suppressed in a sustained reference is exactly the write mechanism here.
\end{remark}

\subsection{Resolution limit}

The recording resolution is bounded by the local lag time through a
time--energy relation: a mark localized to energy width $\Delta E$ cannot be
nucleated faster than $\tau\gtrsim\hbar/\Delta E$, so the product of the
resolvable mark count and the lag time satisfies $\Delta n\cdot\tau\gtrsim
\hbar$. The recorder is thus lag-limited, with the same characteristic time
that governs the propagation of Sec.~\ref{sec:signal} setting its temporal
resolution.

% ============================================================
\section{Threshold Conversion and Reverse-Signature Confirmation}
\label{sec:threshold}
% ============================================================

A leading element delivered to sufficient directed momentum can, on
interaction with a suitable conversion medium, undergo a threshold process: a
portion of its energy is converted into a bound doublet of opposite-coupling
carriers. The process has a sharp energy threshold $E_{\mathrm{th}}$, below
which no doublet forms and above which one may.

\begin{proposition}[Reverse-signature confirmation]
\label{prop:reverse}
Let a forward conversion above threshold produce a bound doublet, and let the
doublet subsequently recombine, emitting two equal quanta in opposite
directions. Then detection of the back-to-back equal-quanta signature confirms
that the forward threshold $E_{\mathrm{th}}$ was crossed, with a
structurally vanishing false-confirmation rate.
\end{proposition}

\begin{proof}
The recombination signature --- two quanta of equal magnitude, emitted
back-to-back and time-coincident --- is a fourfold coincidence. Its occurrence
requires a doublet to have formed, which requires the forward energy to have
exceeded $E_{\mathrm{th}}$. No sub-threshold interaction produces the doublet
and hence none produces the signature. The fourfold coincidence cannot be
assembled from uncorrelated background to any significant rate, so the
confirmation is effectively background-free.\qed
\end{proof}

This realizes detection-by-reverse-process: rather than observe the forward
delivery directly, one observes the unmistakable signature of the reverse
(recombination) process, which exists only if the forward threshold was met.
The threshold makes the confirmation binary; the coincidence makes it pure.
The conversion medium is a thin interaction region, since the doublet forms
only in the presence of a third body to balance momentum.

% ============================================================
\section{Measurement by Timestamped Counting}
\label{sec:counting}
% ============================================================

Given a precise time reference, the measurement of a delivered element reduces
to counting arrivals and recording their times. The directed speed follows
from a single arrival as
\begin{equation}
v = \frac{\Delta x}{\Delta t},
\qquad
\frac{\Delta p}{p} = \frac{\Delta t}{t_{\mathrm{flight}}},
\label{eq:counting}
\end{equation}
where $\Delta x$ is the known flight gap and $\Delta t$ the measured flight
time. No reconstruction of the element's history is required; a timestamped
count is the complete kinematic datum, and the momentum resolution is set
directly by the timing resolution. The recorder of Sec.~\ref{sec:recording}
and the counter of Eq.~\eqref{eq:counting} are complementary readouts of the
same delivery: the former integrates a sequence passively, the latter resolves
single arrivals against a clock.

\begin{remark}
Equation~\eqref{eq:counting} expresses the measurement entirely in terms of
counts and times. The information cost of acquiring a record is the standard
cost of recording a distinguishable outcome \citep{landauer1961}; reading the
fixed record afterward carries no such cost. Acquisition and readout are
thereby cleanly separated.
\end{remark}

% ============================================================
\section{Generality and Instantiation}
\label{sec:discussion}
% ============================================================

\subsection{Substrate independence}

Every result above is a property of the Hamiltonian
Eq.~\eqref{eq:hamiltonian}: a confined ensemble of identical elements with a
repulsive inverse-square coupling. Nothing in the derivations fixes the
physical identity of the elements or the numerical value of $\kappa$. The
signal velocity (Theorem~\ref{thm:disparity}), the compression cascade
(Theorem~\ref{thm:cascade}), the self-confinement (Corollary~\ref{cor:selfconfine}),
the fixed-point selection (Proposition~\ref{prop:detection}), the confirmation
chain (Theorem~\ref{thm:chain}), and the metastable recorder
(Proposition~\ref{prop:recording}) hold for any realization of
Eq.~\eqref{eq:hamiltonian}. Candidate realizations span a wide range of scales,
from charged colloidal ensembles and trapped macroscopic grains to the
one-component plasma of the Wigner lattice \citep{wigner1934,dubin1999,birkl1992}.

\subsection{The regime of instantiation}

The construction acquires its sharpest form when the elements are the
lightest available carriers of the coupling charge and the conversion medium
of Sec.~\ref{sec:threshold} is present. In that regime the chain of results
composes into a single capability: a stochastic source
(Sec.~\ref{sec:release}) feeds a velocity-graded compression cascade
(Sec.~\ref{sec:cascade}) that funnels distributed momentum onto one leading
element over a coupling-set time (Sec.~\ref{sec:signal}); the co-moving
multi-shell lattice self-confines as it speeds up
(Sec.~\ref{sec:transverse}); the coherent delivery is selected by a
fixed-point basin and confirmed by its own existence
(Sec.~\ref{sec:selection}); staged confirmation yields a single, pure outcome
(Sec.~\ref{sec:chains}); a threshold conversion is confirmed by its reverse
signature (Sec.~\ref{sec:threshold}); and the whole measurement reduces to
timestamped counting (Sec.~\ref{sec:counting}), with a metastable variant of
the same lattice available as a passive integrating recorder
(Sec.~\ref{sec:recording}).

The notable feature of this regime is what it does \emph{not} require. It
requires no extended driving structure, because momentum is funnelled at the
signal velocity rather than accumulated along a path. It requires no
high-precision source, because selection is performed downstream by the
fixed-point basin. It requires no reconstruction apparatus, because
confirmation is the existence of the coherent output and measurement is
counting. The spatial extent and control precision that dominate the cost of
conventional directed-delivery systems are, in this construction, replaced by
the intrinsic properties of the coupled lattice. We leave the detailed
engineering of any specific instantiation, and the question of which carriers
and conversion media are most favourable, to subsequent and separate work.

\subsection{The decisive open quantity}

One quantity governs whether the construction functions: the measure
$|\mathcal{B}|$ of the fixed-point basin (Proposition~\ref{prop:detection}),
equivalently the width $\Delta\phi_{\mathrm{acc}}$ of the acceptance window in
release configuration space. If $|\mathcal{B}|>0$ with a computable value,
the yield (Eq.~\eqref{eq:yield}) is nonzero and the detection is binary; if
$|\mathcal{B}|$ is vanishing, no throughput recovers it. Likewise, the
metastable recorder functions only within the stability window of
Definition~\ref{def:metastable}, between self-erasure and global relaxation.
These two windows --- the basin measure and the metastability band --- are the
physical content on which the construction stands or falls, and they are the
proper targets of any test.

% ============================================================
\section{Conclusion}
\label{sec:conclusion}
% ============================================================

We have set out a substrate-independent mechanism, resident in confined
inverse-square-coupled lattices, for delivering a single element to high
directed momentum without active high-precision control, and for recording
sequences of transits passively. The mechanism is organized throughout by a
single inversion (Principle~\ref{prin:inversion}): the mutual repulsion, the
collective instability, and the high failure rate that a control-based design
must suppress are here configured, respectively, as the momentum-funnelling
interaction, the recording mechanism, and the selection filter. The central
enabling fact is the funnelling disparity (Theorem~\ref{thm:disparity}): a
perturbation traverses the coupled lattice far faster than any element moves,
so distributed momentum can be concentrated onto one element over a time set
by the coupling. The central enabling structure is the fixed-point selection
(Proposition~\ref{prop:detection}): the coherent output exists only on a basin
of positive measure, so its existence is its own confirmation. Everything
else --- the graded cascade, the self-confinement, the confirmation chain, the
metastable recorder, the reverse-signature threshold, the counting readout ---
is a consequence or composition of these two facts. The construction's
validity reduces to two measurable quantities: the basin measure and the
metastability band. We have framed the development entirely as properties of a
class of dynamical systems, and we regard the identification of favourable
specific realizations as a separate question, deliberately left open.

% ============================================================
\section*{Acknowledgements}
The author thanks colleagues for discussions on coupled-lattice dynamics and
metastable media.

\bibliographystyle{plainnat}
\bibliography{references}

\end{document}
